\SetPractica{Solubilidad}{Solubilidad}

\section{Introducción teórica}

Las disoluciones son una parte esencial del trabajo de laboratorio, y aunque las propiedades entre diferentes disoluciones varían enormemente, es posible relacionarlas a la naturaleza de sus moléculas y a la forma en que interactúan entre sí.

En primer lugar, las moléculas heteroatómicas cuentan con cierta polaridad, esto a causa de que los electrones de enlace son atraídos con distinta fuerza por cada tipo de átomo. Así, una molécula puede tener cierta polaridad general, a lo cual se le conoce como momento dipolar molecular. Este momento dipolar se relaciona de cierta manera con el concepto de electronegatividad, pues a partir de este se obtienen las diferencias de cargas que indican la polaridad en cada enlace, las cuales se suman para obtener la polaridad general de la molécula. Además, se caracteriza por tener extremos positivos y negativos (Wade, 2011, pp. 9, 10, 61, 63). Este momento dipolar molecular juega un papel fundamental en las interacciones moleculares que veremos a continuación, las cuales se relacionan con la solubilidad entre sustancias. En primer lugar están las fuerzas dipolo-dipolo, las cuales están presentes en moléculas polares, y surgen de la atracción de los extremos positivos y negativos de los momentos dipolares (Wade, 2011, p. 63). En segundo lugar se encuentran las fuerzas de dispersión de London, las cuales se presentan en moléculas no polares y surgen cuando moléculas cercanas inducen momentos dipolares temporales en una molécula. Estos dipolos, aunque son temporales y cambian constantemente, llegan a una fuerza neta de atracción (Wade, 2011, pp. 63-64).

Por último están los puentes de hidrógeno, los cuales se presentan en enlaces OH y NH, donde se deja al átomo de hidrógeno con una carga positiva parcial. Este hidrógeno tiene una gran afinidad electrónica y forma uniones intermoleculares con otros pares de electrones no enlazados. (Wade, 2011, pp. 64-65).

Es sobre esta serie de fuerzas intermoleculares que se forma una famosa regla que dice: “Lo semejante disuelve a lo semejante”. Esta se refiere a que solventes de una naturaleza polar determinada disolverán mejor a solutos de una polaridad semejante. Dado el caso en que tanto solvente como soluto sean polares, las cargas presentes en ambos elementos de la solución permiten que se las uniones se reordenen fácilmente. Por otro lado, si tanto solvente como soluto son no polares, las fuerzas intermoleculares de ambas sustancias son relativamente débiles, por lo que se mezclarán fácilmente, presentando solubilidad. Situaciones de poca o nula solubilidad se darían en los casos en que los componentes de la disolución no compartan polaridad. En el caso de que sea un soluto polar en un solvente no polar, la insolubilidad se daría porque las fuerzas intermoleculares del soluto serían más fuertes que las del solvente. Una situación similar se daría si se intenta disolver un soluto no polar en un solvente polar, pues las fuerzas intermoleculares del solvente serían demasiado fuertes como para que el soluto interaccione con estas y reordene sus uniones, dando paso a la disolución. (Wade, 2011, pp. 66-68)

Además de las interacciones por polaridad, existen otras condiciones que afectan las propiedades de una disolución, como su sobresaturación, caso que se presenta cuando en una solución se excede la cantidad máxima de soluto que es capaz de disolverse en ella. Tal condición se puede identificar por la formación de cristales precipitados en la solución. (CK-12 Foundation, 2019)

\section{Objetivos}

\begin{itemize}
    \item Comprobar de forma experimental como la solubilidad de una sustancia en un determinado disolvente depende de la naturaleza de ambos (soluto y disolvente) y explicar dicho comportamiento mediante las distintas interacciones intermoleculares que puedan establecerse entre las moléculas de soluto y disolvente.
\end{itemize}

\section{Materiales, equipos y reactivos}

\begin{table}[H]
\centering{\small\setstretch{1.25}
\begin{tabular}{|m{0.3\linewidth}|m{0.3\linewidth}|m{0.3\linewidth}|} \hline
    
    \thead{Equipos} & \thead{Materiales} & \thead{Reactivos} \\ \hline
    
    {Balanza digital \par} 
    
    &
    
    {Piseta \par}
    {Tubos de ensayo \par}
    {Gradilla \par}
    {Varilla de agitación \par}
    {Goteros o pipetas Pasteur \par}
    {Pipetas graduadas \par}
    {Perilla \par}
    {Vidrio de reloj \par}
    {Espátula \par}
    {Mrtero y pistilo \par}
    {Embudo}
    
    & 
    
    {Agua destilada \par}
    {Sulfato de cobre (\ce{}) \par}
    {Aceite \par}
    {Parafina líquida \par}
    {Yodo sólido \par}
    {Etanol \par}
    {Hexano \par}
    {Diclorometano}
    
    \\ \hline
\end{tabular}}\end{table}

\section{Procedimientos}

\begin{enumerate}
    \item Poner una pequeña cantidad de sulfato de cobre (0.5g aproximadamente) en tres tubos de ensayo y añadir a continuación: en el primer tubo, 2mL de agua; en el segundo tubo, 2mL de etanol, y en el tercer tubo, 2mL de hexano. Agitar y observar la distinta solubilidad anotando observaciones.
    \item Poner 2 gotas de aceite en tres tubos de ensayo y añadir a continuación: en el primer tubo, 2mL de agua; en el segundo tubo, 2mL de etanol, y en el tercer tubo, 2mL de hexano operando del mismo modo que en el ensayo anterior. 
    \item Poner 2 gotas de parafina líquida en tres tubos de ensayo y añadir a continuación: en el primer tubo, 2mL de agua; en el segundo tubo, 2mL de etanol, y en el tercer tubo, 2mL de hexano operando del mismo modo que en los ensayos anteriores. 
    \item Poner un trocito de yodo en tres tubos de ensayo y añadir a continuación: en el primer tubo, 2mL de agua; en el segundo tubo, 2mL de etanol, y en el tercer tubo, 2mL de diclorometano, operando del mismo modo que en los ensayos anteriores.
    \item Al tubo que contiene diclorometano y yodo, agregar 2mL de agua. Agitar y observar la distribución del soluto entre los dos disolventes inmiscibles entre sí. 
\end{enumerate}